% GENERATED FILE. Edit canonical lessons, then run npm run generate:books.
% canonical-source-hash: fnv1a64:208db4d5c8279b27
% canonical-lessons: LA-C432-equito, LA-C432-ephippium, LA-C432-calcar, LA-C432-desilio, LA-C432-calcitro

\chapter{To ride, Saddle cloth, Spur, To jump down, To kick}
\label{ch:la-to-ride-saddle-cloth-spur-to-jump-down-to-kick}
\hlchaptermodality{432}

\begin{chapteropening}
\textbf{By the end of this chapter:} \emph{I can talk about riding: riding a horse, the saddle cloth, a spur, jumping down, and a horse kicking.}

Complete the last lesson of chapter 432: I can talk about riding: riding a horse, the saddle cloth, a spur, jumping down, and a horse kicking.
\end{chapteropening}

\section[equitō, equitāre]{equitō, equitāre — to ride (on horseback)}

\label{lesson:LA-C432-equito}



\begin{quote}
Before the new one: say the Latin for to bridle, then the Latin for a whip.
\end{quote}

\subsection*{You'll want to know: equitō, equitāre}
\textbf{equitō, equitāre} — \textquotedblleft{}to ride (on horseback)\textquotedblright{}.

\begin{quote}
Mīles per campum equitat.
\end{quote}

\begin{quote}
\emph{The soldier rides across the plain.}
\end{quote}

\emph{equitō, equitāre, equitāvī, equitātum}. English \emph{equitation} comes from it.

\begin{cousinweb}
\emph{equitō} is built on \emph{eques} 'horseman' (EtymDB; Lewis \& Short), which comes from \emph{equus} 'horse'. English \emph{equitation} comes from it.
\end{cousinweb}

\begin{grammarlens}[title={What you've built}]
One more word for telling, explaining and describing at B1.
\end{grammarlens}

\subsection*{Your turn}
\begin{itemize}
\raggedright
  \item \emph{Say it:} \emph{equitō, equitāre}
  \item \emph{Say it:} \emph{equitō, equitāre}, once more
  \item \emph{Say it:} a sentence of your own with \emph{equitō, equitāre}
  \item \emph{Recall:} say the Latin for to bridle, then the Latin for a whip, then say \emph{equitō, equitāre} again
\end{itemize}

\begin{tcolorbox}[breakable,colback=teal!4,colframe=teal!35!black,title={Before you move on}]
What does equitō, equitāre mean? (\textquotedblleft{}To ride (on horseback)\textquotedblright{}.) Say it once more.
\end{tcolorbox}
\section[ephippium, ephippiī]{ephippium, ephippiī — a saddle cloth, a horse blanket}

\label{lesson:LA-C432-ephippium}



\begin{quote}
Before the new one: say the Latin for a whip, then the Latin for to ride.
\end{quote}

\subsection*{You'll want to know: ephippium, ephippiī}
\textbf{ephippium, ephippiī} — \textquotedblleft{}a saddle cloth, a horse blanket\textquotedblright{}.

\begin{quote}
Servus ephippium in equum pōnit.
\end{quote}

\begin{quote}
\emph{The slave puts the saddle cloth on the horse.}
\end{quote}

Second declension, neuter: \emph{ephippium, ephippiī n.} Riders sat on a padded cloth rather than a modern saddle.

\begin{cousinweb}
\emph{ephippium} was borrowed from Greek \emph{ephippion} 'saddlecloth', built from \emph{epi} 'on' and \emph{hippos} 'horse' (EtymDB; Lewis \& Short).
\end{cousinweb}

\begin{grammarlens}[title={What you've built}]
One more word for telling, explaining and describing at B1.
\end{grammarlens}

\subsection*{Your turn}
\begin{itemize}
\raggedright
  \item \emph{Say it:} \emph{ephippium, ephippiī}
  \item \emph{Say it:} \emph{ephippium, ephippiī}, once more
  \item \emph{Say it:} a sentence of your own with \emph{ephippium, ephippiī}
  \item \emph{Recall:} say the Latin for a whip, then the Latin for to ride, then say \emph{ephippium, ephippiī} again
\end{itemize}

\begin{tcolorbox}[breakable,colback=teal!4,colframe=teal!35!black,title={Before you move on}]
What does ephippium, ephippiī mean? (\textquotedblleft{}A saddle cloth, a horse blanket\textquotedblright{}.) Say it once more.
\end{tcolorbox}
\section[calcar, calcāris]{calcar, calcāris — a spur}

\label{lesson:LA-C432-calcar}



\begin{quote}
Before the new one: say the Latin for to ride, then the Latin for a saddle cloth.
\end{quote}

\subsection*{You'll want to know: calcar, calcāris}
\textbf{calcar, calcāris} — \textquotedblleft{}a spur\textquotedblright{}.

\begin{quote}
Eques calcar habet, sed equus lentē it.
\end{quote}

\begin{quote}
\emph{The rider has a spur, but the horse goes slowly.}
\end{quote}

Third declension, neuter: \emph{calcar, calcāris n.} Figurative too: a spur to effort.

\begin{cousinweb}
\emph{calcar} is usually linked with \emph{calx, calcis} 'heel': Lewis \& Short compare \emph{calx} and quote an ancient explanation that the spur is tied to the heel. EtymDB lists the word but gives no parent.
\end{cousinweb}

\begin{grammarlens}[title={What you've built}]
One more word for telling, explaining and describing at B1.
\end{grammarlens}

\subsection*{Your turn}
\begin{itemize}
\raggedright
  \item \emph{Say it:} \emph{calcar, calcāris}
  \item \emph{Say it:} \emph{calcar, calcāris}, once more
  \item \emph{Say it:} a sentence of your own with \emph{calcar, calcāris}
  \item \emph{Recall:} say the Latin for to ride, then the Latin for a saddle cloth, then say \emph{calcar, calcāris} again
\end{itemize}

\begin{tcolorbox}[breakable,colback=teal!4,colframe=teal!35!black,title={Before you move on}]
What does calcar, calcāris mean? (\textquotedblleft{}A spur\textquotedblright{}.) Say it once more.
\end{tcolorbox}
\section[dēsiliō, dēsilīre]{dēsiliō, dēsilīre — to jump down, to leap down}

\label{lesson:LA-C432-desilio}



\begin{quote}
Before the new one: say the Latin for a saddle cloth, then the Latin for a spur.
\end{quote}

\subsection*{You'll want to know: dēsiliō, dēsilīre}
\textbf{dēsiliō, dēsilīre} — \textquotedblleft{}to jump down, to leap down\textquotedblright{}.

\begin{quote}
Ad portam eques ex equō dēsilit.
\end{quote}

\begin{quote}
\emph{At the gate the rider jumps down from his horse.}
\end{quote}

\emph{dēsiliō, dēsilīre, dēsiluī, dēsultum}, with \emph{ex} or \emph{dē}. English \emph{desultory} comes from \emph{dēsultor}, a circus rider who leapt from horse to horse.

\begin{cousinweb}
\emph{dēsiliō} is built from \emph{dē-} 'down' and \emph{saliō} 'to leap, to jump' (EtymDB; Lewis \& Short). English \emph{desultory} comes from \emph{dēsultōrius}, built on its participle.
\end{cousinweb}

\begin{grammarlens}[title={What you've built}]
One more word for telling, explaining and describing at B1.
\end{grammarlens}

\subsection*{Your turn}
\begin{itemize}
\raggedright
  \item \emph{Say it:} \emph{dēsiliō, dēsilīre}
  \item \emph{Say it:} \emph{dēsiliō, dēsilīre}, once more
  \item \emph{Say it:} a sentence of your own with \emph{dēsiliō, dēsilīre}
  \item \emph{Recall:} say the Latin for a saddle cloth, then the Latin for a spur, then say \emph{dēsiliō, dēsilīre} again
\end{itemize}

\begin{tcolorbox}[breakable,colback=teal!4,colframe=teal!35!black,title={Before you move on}]
What does dēsiliō, dēsilīre mean? (\textquotedblleft{}To jump down, to leap down\textquotedblright{}.) Say it once more.
\end{tcolorbox}
\section[calcitrō, calcitrāre]{calcitrō, calcitrāre — to kick (with the heels); to resist}

\label{lesson:LA-C432-calcitro}



\begin{quote}
Before the new one: say the Latin for a spur, then the Latin for to jump down.
\end{quote}

\subsection*{You'll want to know: calcitrō, calcitrāre}
\textbf{calcitrō, calcitrāre} — \textquotedblleft{}to kick (with the heels); to resist\textquotedblright{}.

\begin{quote}
Mūlus calcitrat, et servus fugit.
\end{quote}

\begin{quote}
\emph{The mule kicks, and the slave runs away.}
\end{quote}

A first-conjugation verb; the perfect is rare. English \emph{recalcitrant} comes from \emph{recalcitrāre} 'to kick back'.

\begin{cousinweb}
\emph{calcitrō} is built on \emph{calx, calcis} 'heel' (EtymDB; Lewis \& Short): to strike out with the heels. English \emph{recalcitrant} comes from the compound \emph{recalcitrāre}.
\end{cousinweb}

\begin{grammarlens}[title={What you've built}]
One more word for telling, explaining and describing at B1.
\end{grammarlens}

\subsection*{Your turn}
\begin{itemize}
\raggedright
  \item \emph{Say it:} \emph{calcitrō, calcitrāre}
  \item \emph{Say it:} \emph{calcitrō, calcitrāre}, once more
  \item \emph{Say it:} a sentence of your own with \emph{calcitrō, calcitrāre}
  \item \emph{Recall:} say the Latin for a spur, then the Latin for to jump down, then say \emph{calcitrō, calcitrāre} again
\end{itemize}

\begin{tcolorbox}[breakable,colback=teal!4,colframe=teal!35!black,title={Before you move on}]
What does calcitrō, calcitrāre mean? (\textquotedblleft{}To kick (with the heels); to resist\textquotedblright{}.) Say it once more.
\end{tcolorbox}
